\subsubsection{What canonical positivity does and does not assert}

The strict effectivity assumptions in \fid{F57}--\fid{F58}
are additional comparison hypotheses.
They must not be substituted for the published
signed/factored positivity statements.

VIII, Conjecture~3.4, \S3.5, p.~24, allows each relevant
right-algebra structure coefficient to have the form
\[
       \pm(q^{1/2}-q^{-1/2})^a f(q),
       \qquad a\in\mathbb Z_{\ge0},
\]
with a bar-invariant positive unimodal factor \(f\).
Positive there means nonnegative rational coefficients.
The clause includes internal multiplication constants
and the external expansion coefficients in equation~(13).
It applies to both \(E(r)\) and \(E^{\mathrm{opp}}(r)\)
\cite{VIII}.
The sign and vanishing factor are part of the stated
conjecture.

This distinction is already visible in the published
rank-three example. Figure~2, p.~34, gives
\begin{equation}\label{eq:publishedsignedcanonicalconstant}
 c(\gamma_2^1,\gamma_1^1;\gamma_{21}^1)
       =-\frac{(q^2+q+1)(q-1)^2}{q^2}.
\end{equation}
Its sign is negative for positive \(q\ne1\).
With \(v=q^{1/2}\), it is
\[
              -(v-v^{-1})^2(v^2+1+v^{-2}).
\]
The source presents these constants as evidence for
Conjecture~3.4; see \S5.1.2, p.~29.
This is a cited source computation, not a newly
recomputed full multiplication table.
The displayed factorization is immediate.

The allowed signed/factored class is not closed under
addition. For example \(1\) and \(-(v-v^{-1})^2\)
belong to it, whereas their sum
\(3-v^2-v^{-2}\) does not.
Its nonzero value at \(v=1\) rules out a positive
vanishing-factor exponent, and its mixed signs rule
out exponent zero with either overall sign.
A statement about individual canonical constants
therefore does not give the same statement for every
named-word trace. The \fid{F57} trace likewise has
value \(4096\) at \(v=1\) and mixed coefficient signs,
but it is not an individual constant covered by
Conjecture~3.4.

The opening of VIII, \S3, p.~18, calls the named
\(\mathcal P_i,\mathcal Q_i\) positive and negative
generators because they rescale the corresponding
spectral projectors from VII, \S5.
Their membership in the ambient integral form and their
ring-theoretic generation of the algebra do not assert
membership in an effective canonical cone.
The minimum-degree construction in Conjectures~3.1--3.2,
pp.~21--23, is a different condition.
Its definition does not identify every named generator
with a canonical basis element or an effective Laurent
combination of such elements \cite{VII,VIII}.

The unsigned external-expansion clause of VIII,
Conjecture~2.7, p.~15, concerns \(CB(\mathcal A_r^H)\)
on the nonstandard Schur side, in the tensor basis
specified before equation~(9).
The named \(P_i\) belong to the commuting right algebra
\(\mathcal B_r^H\). No identification with those left
canonical elements is supplied by that clause.
Even there the internal multiplicative coefficients
retain signs and vanishing factors.

On the revised IV side, Definition~5.3, p.~29, imposes
weight, lattice, bar, and crystal conditions for an
ordinary upper based module; Definition~13.2, p.~73,
allows the additional swap to carry a basis vector to
either sign of another.
Conjecture~19.1, pp.~130--131, asks for the specified
generic cells, specialized quasi-cells and fat-cell
subquotients. It contains no unsigned \(P_i\)- or
\(Q_i\)-matrix condition.
Remark~11.19, p.~67, and Example~19.3, pp.~133--135,
explicitly retain signed/factored coefficients,
including \(c_0=(v-v^{-1})^2\).
The discussion following Conjecture~19.1 says that a
based-module theory for the nonstandard Schur algebra
would be an additional strengthening \cite{IV}.

\paragraph{A precise stronger KL-style package.}
Let \(R=\mathbb Q[v,v^{-1}]\) and
\(R_+=\mathbb Q_{\ge0}[v,v^{-1}]\).
Suppose an \(R\)-algebra \(A\) is free on a finite basis
\(C\), every product of two elements of \(C\) has
coefficients in \(R_+\), and every named \(P_i\) lies
in \(R_+C\).
If a module \(M\) is a coordinate cellular subquotient
of the regular module with its \(C\) basis, then every
\(P_i\) acts on its quotient basis with coefficients
in \(R_+\).

Indeed write \(P_i=\sum_c h_{i,c}c\), with
\(h_{i,c}\in R_+\). Its multiplication matrix is a
nonnegative Laurent combination of the multiplication
matrices of the \(c\)'s.
Restriction to a stable basis-subset span and passage
to a quotient by another such span retain the
corresponding matrix entries, without subtraction.
They are still in \(R_+\).

Thus unsigned basis multiplication, effective membership
of the named generators, and the coordinate-cell
comparison are separate premises.
If such cells realize all actual generic simples of the
full rank-four algebra, they include \(N^\#\).
The resulting positive \(P_i\) matrices contradict
\fid{F57}.
The weaker direct action premise
\[
                       P_iC_M\subseteq R_+C_M
\]
already suffices for that contradiction.
It is not a consequence of the signed
Conjecture~3.4 as printed.

For the two-row tensor basis, adding this unsigned
action premise to the exact generic and specialized
cell conditions of IV, Conjecture~19.1(ii)--(iii),
instead gives the contradiction in \fid{F58}.
VIII, Conjecture~3.7, pp.~25--26, expressly allows
the necessary change of zero pattern: equation~(15)
defines quasi-subcells using the terms whose vanishing
exponent is zero, and terms of positive exponent
disappear at \(q=1\).
Identifying its opposite algebra basis with the
revised IV tensor basis would itself require a
further comparison \cite{VIII,IV}.

\paragraph{Two exact normalization limits.}
The negative trace also survives any replacement
\[
 P_i'=\epsilon_i h_i(v)P_i,\qquad
 \epsilon_i\in\{1,-1\},\quad
 0\ne h_i\in\mathbb Q_{\ge0}[v,v^{-1}].
\]
For the specified \fid{F57} word its trace becomes
\(\epsilon H(v)p((v+v^{-1})^2)\), where \(H\) is a
nonzero nonnegative Laurent polynomial and
\(\epsilon=\pm1\).
If \(\epsilon=1\), its highest Laurent coefficient
is negative because that of \(p\) is \(-1\).
If \(\epsilon=-1\), its value at \(v=1\) is
\(-4096H(1)<0\), so some coefficient is negative.
These signed positive-Laurent generator rescalings
therefore cannot produce strict positive matrices
on the specified summand.

In contrast, dividing out \(v-v^{-1}\) is not a
regular normalization at the identity fiber.
It is a nonunit of
\(\mathbb Q[v,v^{-1}]_{(v-1)}\).
After inverting it no specialization sending
\(v\) to \(1\) exists, since it would send a unit
to zero. Diagonal basis changes by units of this
local ring preserve the order-two vanishing of
the \(c_0\) edge. General basis mixing must be checked
against the designated coordinate-submodule premise;
if that premise
is preserved, \fid{F58} still applies.
Coefficientwise removal of factors therefore does
not by itself construct coherent relation maps or
preserve the designated specialization.

These conclusions leave the normalized positivity
and unimodality questions in their stated scope.
Mixed or nonproper realizations may account for raw
signed coefficients; the torus computation above
does so for \(c_0\).
An effective graded dimension, a pure point-complex
Euler trace, a pure coefficient on a source, and an
inserted-operator trace are distinct comparisons.

% Separate paragraph for Part III.
\paragraph{The named-generator gate in Grothendieck's language.}
The effective representation cone of a split torus is
closed under addition and tensor product, whereas its
Grothendieck group also admits subtraction.
SGA~3, I, 4.7.2--3 and cote~90, scan~17 supply the
action framework; SGA~6, V, \S\S2--3 and cote~133,
scan~18 supply the coefficient representation
vocabulary \cite{SGA3,SGA6,Archive}.
An algebra element is not automatically an effective
object in that cone. Identifying the named \(P_i\)
with such objects and their cell actions with positive
decomposition multiplicities are precisely the extra
hypotheses proved sufficient above.
The signed/factored source constants live naturally
as virtual characters or require an additional
normalization. This framework imposes no
\(\lambda\)-ring structure on the noncommutative word
algebra and supplies no unproved canonical-basis
comparison.
